% ══════════════════════════════════════════════════════════════════════════════
%  EXPERIMENT 8 — DISTILLATION                                   (Lesson 17, Term 2 Week 6)
% ══════════════════════════════════════════════════════════════════════════════
\experimentdiv{Distillation}{Setting up a distillation apparatus \textbullet{} Heating a liquid safely \textbullet{} Testing the distillate}{Lesson 17 \textbullet\ Term 2, Week 6 \textbullet\ Lab}\label{sec:8}
\lessoncodes{VC2S8U05, U07 \textbullet{} VC2S8I02, I03, I06, I08}

\vocabhead
\begin{ctab}
\begin{tabular}{|SW{4.2cm}|W{11.6cm}|}
  \hline\hd & \thead{Word} & \thead{Definition (copy from the slide)} \\ \hline
  & \scaf{\textbf{Distillation}} & \scaf{Separating a liquid from a solution by boiling it and condensing the vapour.} \\ \hline
  & \textbf{Distillate} & \\ \hline
  & \textbf{Condense} & \\ \hline
  & \textbf{Boiling point} & \\ \hline
  & \textbf{Boiling chip} & \\ \hline
\end{tabular}
\end{ctab}

\begin{lipanel}
\inquiryli{Select and set up equipment for distillation, and use it safely.}
\understandli{Explain how distillation separates pure water from a solution, using changes of state.}
\sctitle
\begin{itemize}[leftmargin=1.4em, itemsep=0pt, topsep=1pt]
  \item I can name each piece of equipment and say what it does.
  \item I can draw and label the distillation set-up.
  \item I can explain where the water went and what was left behind.
\end{itemize}
\end{lipanel}

\begin{skillsfocus}
\focus{PC7}{Choose and set up the distillation equipment.}
\focus{PC3}{Heat a flask safely.}
\focus{PC6}{Test the distillate and the residue.}
\focus{CO2}{Draw the set-up as a scientific diagram.}
\focus{EV2}{Compare the results with the water cycle.}
\end{skillsfocus}

% ─── TECHNIQUE ────────────────────────────────────────────────────────────────
\exppart{Technique spotlight}{Blue}

\begin{technique}{The distillation set-up}{PC7}
\begin{center}
\begin{tikzpicture}[scale=0.78]
  % stand and clamp holding the flask
  \begin{scope}[shift={(-1.6,0)}]\draw[eq] (0,0) rectangle (2.0,0.2); \draw[eq] (0.3,0.2) -- (0.3,5.6) (0.45,0.2) -- (0.45,5.6);\end{scope}
  \draw[eq] (-1.15,4.6) -- (0.6,4.6) (-1.15,4.75) -- (0.6,4.75);
  % tripod, gauze, Bunsen
  \begin{scope}[shift={(-0.35,0.2)}]\eqBunsen\end{scope}
  \begin{scope}[shift={(-0.6,2.2)}]\draw[eq] (-0.15,0.0) rectangle (2.55,0.1);\draw[eq] (0.25,0) -- (0,-1.95) (2.15,0) -- (2.4,-1.95);\end{scope}
  % conical flask
  \begin{scope}[shift={(-0.15,2.3)}]\eqFlask\end{scope}
  \fill[black!15] (-0.12,2.33) -- (0.22,3.3) -- (1.53,3.3) -- (1.72,2.33) -- cycle;
  \draw[Ink, line width=0.8pt] (0.22,3.3) -- (1.53,3.3);
  \foreach \p in {(0.5,2.5),(0.9,2.6),(1.2,2.45)}{\draw[Ink, fill=white] \p circle (0.07);}
  \fill[Ink] (0.45,5.1) rectangle (1.15,5.45);
  % delivery tube
  \draw[Ink, line width=1.6pt] (0.8,5.45) -- (0.8,5.8) -- (4.4,4.3) -- (4.4,2.0);
  % ice bath and collecting tube
  \begin{scope}[shift={(3.5,0.2)}]\eqBeaker\end{scope}
  \fill[black!6] (3.52,0.22) rectangle (5.28,2.0); \draw[Ink, line width=0.8pt] (3.52,2.0) -- (5.28,2.0);
  \foreach \p in {(3.8,1.6),(5.0,1.2),(3.9,0.7),(5.0,0.5)}{\draw[Ink] \p ++(-0.15,-0.15) rectangle ++(0.3,0.3);}
  \begin{scope}[shift={(4.1,0.4)}]\eqTestTube\end{scope}
  \fill[black!3] (4.12,0.7) arc[start angle=180,end angle=360,radius=0.28] -- (4.68,1.0) -- (4.12,1.0) -- cycle;
  \eqlabl{(0.2,3.0)}{(-2.2,3.6)}{salty coloured water}
  \eqlabl{(0.5,2.5)}{(-2.2,2.6)}{boiling chips}
  \eqlabl{(0.45,5.3)}{(-2.2,5.9)}{rubber bung}
  \eqlab{(2.6,5.05)}{(3.6,5.9)}{delivery tube}
  \eqlab{(5.28,1.5)}{(6.0,1.9)}{ice and water}
  \eqlab{(4.68,0.85)}{(6.0,0.9)}{distillate}
  \eqlabl{(-0.1,0.8)}{(-2.2,1.4)}{Bunsen burner}
\end{tikzpicture}
\end{center}
\begin{itemize}[leftmargin=1.4em, itemsep=1pt, topsep=2pt]
  \item The \textbf{delivery tube} slopes \emph{down} to the collecting tube, so the distillate runs away from the flask.
  \item The collecting tube sits in \textbf{ice water}, which cools the vapour so that it condenses.
  \item \textbf{Boiling chips} make the liquid boil smoothly instead of in sudden violent bursts.
  \item The end of the delivery tube stays \textbf{above} the liquid in the collecting tube, or cold water can be sucked back into the hot flask.
\end{itemize}
\end{technique}

\begin{youdobox}[Choosing equipment]
\task Complete the table. \skr{PC7}
\begin{ctab}
\begin{tabular}{|SW{4.0cm}|W{11.0cm}|}
  \hline\hd & \thead{Equipment} & \thead{Its job in distillation} \\ \hline
  & \scaf{Bunsen burner} & \scaf{Heats the solution until the water boils} \\ \hline
  & Conical flask & \\ \hline
  & Rubber bung & \\ \hline
  & Delivery tube & \\ \hline
  & Beaker of ice water & \\ \hline
  & Boiling chips & \\ \hline
\end{tabular}
\end{ctab}
\end{youdobox}

% ─── PLAN AND SAFETY / METHOD ────────────────────────────────────────────────
\exppart{Plan, safety and method}{Teal}

\textbf{Aim:} To separate pure water from salty, coloured water by distillation.

\task Predict the colour of the distillate and what will be left in the flask. Explain your prediction. \skr{QP1}
\answerbox{4cm}

\subsubsection*{Equipment}
250\,mL conical flask \textbullet{} rubber bung with a bent delivery tube \textbullet{} retort stand and clamp \textbullet{} tripod, gauze mat, Bunsen burner and heatproof mat \textbullet{} test tube \textbullet{} 250\,mL beaker of ice and water \textbullet{} 3 boiling chips \textbullet{} 50\,mL of salt water with blue food colouring \textbullet{} 2 watch glasses \textbullet{} dropper

\begin{infobox}{Safety}
Wear safety glasses. Your teacher checks the set-up before you light the Bunsen. Never heat a sealed container: the delivery tube must be open at the end. Turn the gas off \emph{before} the flask boils dry, and lift the delivery tube out of the collecting tube when you stop heating. The flask, tube and steam stay hot for several minutes.
\end{infobox}

\subsubsection*{Method}
\begin{enumerate}[itemsep=2pt]
  \item Put 50\,mL of salty coloured water and three boiling chips into the conical flask.
  \item Clamp the flask above the tripod and gauze, and fit the bung and delivery tube.
  \item Stand the test tube in the ice water with the end of the delivery tube inside it, above the bottom.
  \item Ask your teacher to check the set-up, then heat gently with a blue flame.
  \item When the test tube is about a quarter full, lift the delivery tube out, then turn off the gas.
  \item Put one drop of the distillate on a watch glass and one drop of the original solution on another. Leave them on the windowsill until next lesson, then compare.
\end{enumerate}

% ─── RESULTS ─────────────────────────────────────────────────────────────────
\exppart{Results}{Amber}

\subsubsection*{Table 8.1: Observations during distillation \skr{PC6}}
\begin{ctab}
\begin{tabular}{|SW{5.0cm}|W{10.0cm}|}
  \hline\hd & \thead{What was observed} & \thead{Observation} \\ \hline
  & \scaf{Solution in the flask at the start} & \scaf{Clear, blue liquid} \\ \hline
  & Inside the flask while it boiled & \\ \hline
  & Inside the delivery tube & \\ \hline
  & Distillate in the test tube & \\ \hline
  & Liquid left in the flask & \\ \hline
  & Watch glass: drop of distillate (next lesson) & \\ \hline
  & Watch glass: drop of original solution & \\ \hline
\end{tabular}
\end{ctab}

\needspace{20\baselineskip}
\task Draw a scientific diagram of your set-up without looking at the technique box. Label every piece of equipment, the solution and the distillate. \skr{CO2}
\drawbox{9cm}{}

% ─── ANALYSIS ────────────────────────────────────────────────────────────────
\exppart{Analysis and evaluation}{Coral}

\task Describe what happened to the water, naming the two changes of state.
\answerbox{4cm}

\needspace{10\baselineskip}
\task How do the watch glasses show that the distillate was pure water? \skr{PC6}
\answerbox{4cm}

\needspace{10\baselineskip}
\task Why did the salt and colouring stay in the flask? Use the idea of boiling point.
\answerbox{4cm}

\needspace{10\baselineskip}
\task Compare distillation with the water cycle. Which part of the water cycle is like the Bunsen burner, and which is like the ice water? \skr{EV2}
\answerbox{4.5cm}

\needspace{10\baselineskip}
\task Never taste the distillate, even though it should be pure water. Give two reasons.
\answerbox{3.5cm}

\begin{yr8box}
\begin{keepext}
\exttask Wonthaggi's desalination plant uses reverse osmosis rather than distillation to make drinking water from seawater. Suggest why distillation would cost much more to run for a whole city.
\answerbox{4.5cm}
\end{keepext}
\end{yr8box}

% ─── SKILL CHECK ─────────────────────────────────────────────────────────────
\exppart{Skill check}{NavyBlue}

\begin{skillcheck}
\scrow{PC7}{I can name each piece of distillation equipment and its job.}
\scrow{PC3}{I heated the flask safely and stopped before it boiled dry.}
\scrow{PC6}{I tested the distillate and compared it with the solution.}
\scrow{CO2}{I can draw and label the set-up from memory.}
\scrow{EV2}{I can compare distillation with the water cycle.}
\end{skillcheck}

\task Which skill will you focus on in Experiment 9? Explain what you will do differently.
\answerbox{5cm}
